% Generated draft: requires mathematical and source review.
\documentclass[11pt]{article}
\usepackage[margin=25mm]{geometry}
\usepackage{fontspec}
\setmainfont{Latin Modern Roman}
\usepackage{amsmath,amssymb,mathtools,unicode-math}
\setmathfont{Latin Modern Math}
\usepackage{enumitem,needspace,xurl,hyperref}
\setlength{\parindent}{0pt}
\setlength{\parskip}{5pt}
\setlength{\emergencystretch}{3em}
\newcommand{\N}{\mathbb{N}}
\newcommand{\Z}{\mathbb{Z}}
\newcommand{\Q}{\mathbb{Q}}
\newcommand{\R}{\mathbb{R}}
\newcommand{\C}{\mathbb{C}}
\newcommand{\F}{\mathbb{F}}
\newcommand{\A}{\mathbb{A}}
\newcommand{\PP}{\mathbb{P}}
\newcommand{\E}{\mathbb{E}}
\newcommand{\eps}{\varepsilon}
\newcommand{\dd}{\,\mathrm{d}}
\begin{document}
\section*{UnsolvedMath ID: 30004225}
\textbf{Source problem code:} catalogue entry
\subsection{Definitions and mathematical statement}
The supplied record asks two questions about polytope face vectors. Its notation is not fully defined, so the following separates the recorded questions from conditional interpretations.

Under the usual interpretation, \(\mathcal P^d\) denotes the class of convex \(d\)-dimensional polytopes, \(\sigma^d\) denotes a \(d\)-simplex, and \(f_i(P)\) counts the \(i\)-dimensional faces of \(P\). The record does not explicitly specify this class or the coordinates included in \(f(P)\).

\textbf{Limit-ray question.} For every integer \(d>4\), is every ray of
\[
\operatorname{Cone}_{f(\sigma^d)}\bigl(f(\mathcal P^d)\bigr)
\]
a limit ray? Here \(f(\mathcal P^d)=\{f(P):P\in\mathcal P^d\}\). The supplied record does not define the cone construction, whether closure is taken, the placement of its apex, or what qualifies as a limit ray. In particular, it does not specify the sequences, normalization, topology, or any unboundedness requirement used in defining a limit ray. These conventions must be recovered before this question has a precise mathematical formulation.

\textbf{Four-dimensional ratio question, as recorded.} Does there exist a real constant \(C>0\), independent of \(P\), such that for every \(P\in\mathcal P^4\),
\[
\frac{f_1(P)+f_4(P)}{f_0(P)+f_3(P)}\le C?
\]
The index \(4\) in the numerator is retained exactly from the source record. With ordinary face-dimension indexing, a four-dimensional polytope has \(f_4(P)=1\) if the whole polytope is counted as a face; if its face vector contains only proper faces, \(f_4(P)\) is not a coordinate. Under the former convention and the convex-polytope interpretation, the denominator is positive and the displayed question becomes
\[
\exists C>0\ \forall P\in\mathcal P^4:\qquad
f_1(P)+1\le C\bigl(f_0(P)+f_3(P)\bigr).
\]
Whether this literal expression is the intended fatness parameter is unresolved. No replacement of \(f_4\) or alternative indexing convention is assumed.
\subsection{Short English statement}
The record asks whether every ray of a simplex-based face-vector cone is a limit ray in each dimension \(d>4\), and whether the displayed four-dimensional ratio admits a uniform upper bound. The cone and limit-ray conventions are missing, and the numerator's \(f_4\) requires verification. Both questions are supplied as open; this draft does not independently verify that status.
\subsection{Sources}
\begin{itemize}
\item Local UnsolvedMath registry, permanent ID \(30004225\), source problem code \texttt{catalogue entry}, title: \emph{Limit Rays and Fatness of Polytope Face Vector Cones}. The supplied record belongs to the Oberwolfach Reports Open Problems collection, lists proposed year 2019, and names no proposer.
\item Supplied report citation: \emph{Geometric, Algebraic, and Topological Combinatorics} (2020), Oberwolfach Reports, pp.~2395--2472, DOI \texttt{10.4171/owr/2019/39}. \url{https://doi.org/10.4171/owr/2019/39}. The background describes these reports as documenting mathematical workshops and their open-problem discussions. This is the supplied authoritative provenance; the report text has not been independently consulted here.
\item Supplied literature source: \emph{Projected products and fatness}. \url{https://arxiv.org/abs/math/0407042}. No authors or specific supporting results are supplied in the record.
\item Literature assessment dated \texttt{2026-08-22}: the supplied curation labels the problem open and states that known \(d=4\) ray-limit and fatness bounds do not settle the requested cone description or finite-generation question, while the corresponding higher-dimensional questions remain open. This is supplied dated curation, not an independently verified literature assessment. The background explicitly requests independent verification of time-sensitive status and cited claims. Its mention of finite generation is not accompanied by a mathematical formulation in the recorded questions.
\item The registry labels the statement \texttt{corrected\_verified} and says that mathematical content was preserved while notation, typography, and sentence structure were repaired. This records a supplied editorial assessment and does not resolve the formulation issues identified here.
\end{itemize}
\end{document}
